\scanpage{93}
\mspage{46}
où $I_{s'}$ \deleted{est le groupe} \inserted{$\simeq\mathcal E_{\widetilde K}$}, $\widetilde{\widetilde K}$ étant le corps des fractions du hensélisé (\deleted{\illegible{2}} strict) de $\underline{\mathcal O}_{S,s}$. ($N_{s'}$ \deleted{est celui du corps de} \inserted{\unsure{s'identifie de|s'identifie à} $\unsure{=}$ à} $\mathcal E_{\widetilde K}$, où $\widetilde K$ correspond au hensélisé pas strict)

Ceci posé, on trouve comme précédemment, si le composé \deleted{de} $\mathcal E_K\mathrel{\unsure{\longrightarrow|-}}\mathcal E_{U_\eta}\mathrel{\unsure{\longrightarrow|-}}\mathcal E_U$ s'annule sur $I_{s'}$ \inserted{(a priori, il envoie $I_{s'}$ dans $\unsure{\Pi|T}$)} \deleted{\unsure{pour}} (ce qui signifie justement que $\unsure{s'|s}\in S'\unsure{\ldots}$) \deleted{que} \unsure{un hom|on trouve} $\mathcal E_{k(s)}\longrightarrow\mathcal E_U$, ou si on préfère $\mathcal E_{k(s)}\longrightarrow\mathcal E_{U_{\widetilde{\mathcal O}_s}}\xleftarrow{\sim}\mathcal E_{X_s}$, qui \deleted{peut} \inserted{donc} s'inter\unsure{prète|préter} comme une section de $\mathcal E_{X_s}$ sur $\mathcal E_{k(s)}$. Et que se passe-t-il pour $s\notin S'$ ? \deleted{\unsure{Alors}} On a vu que $D=S\setminus S'$ est purement de codim $1$ (c'est l'image réciproque du diviseur $T$ sur $X$ par la fonction $f$ de $X$ sur $S$, définie si $g\geq1$ ou si $S$ de dim $1\ldots$) en fait le $\unsure{\mu=\mu_\alpha|n=n_\alpha}$ de \unsure{tantôt}, relatif à
\margin{92}
